\section{Поточечная и равномерная сходимость}

Пусть на множестве $X\subset\R$ задана последовательность функций $\{f_n(x)\}$.

\begin{definition}[поточечная сходимость]
$f_n\to f$ \emph{поточечно} на $X$, если при каждом фиксированном $x\in X$ числовая
последовательность $\{f_n(x)\}$ сходится к $f(x)$:
\[
  \forall x\in X\ \forall\eps>0\ \exists N(\eps,x)\ \forall n>N:\ \abs{f_n(x)-f(x)}<\eps .
\]
Здесь номер $N$ зависит и от $x$.
\end{definition}

\begin{definition}[равномерная сходимость]
$f_n\rightrightarrows f$ \emph{равномерно} на $X$, если номер $N$ можно выбрать общим для всех $x$:
\[
  \forall\eps>0\ \exists N(\eps)\ \forall n>N\ \forall x\in X:\ \abs{f_n(x)-f(x)}<\eps .
\]
Геометрически: начиная с номера $N$, весь график $f_n$ лежит в $\eps$-полосе вокруг графика $f$.
\end{definition}

\begin{criterion}[супремум-критерий]
$f_n\rightrightarrows f$ на $X$ тогда и только тогда, когда
$\displaystyle\rho_n=\sup_{x\in X}\abs{f_n(x)-f(x)}\to0$.
\end{criterion}

\begin{example}\leavevmode
\begin{enumerate}[label=\arabic*),leftmargin=*]
  \item $f_n(x)=x^n$ на $[0,1)$: поточечно $f_n\to0$, но
        $\rho_n=\sup_{[0,1)}\abs{x^n}=1\not\to0$ — сходимость не равномерна.
  \item «Бегущий горб» $f_n(x)=nxe^{-nx^2}$ на $[0,1]$: поточечно $f_n\to0$, однако в точке
        $x_n=\tfrac{1}{\sqrt{2n}}$ имеем $f_n(x_n)=\sqrt{\tfrac{n}{2}}\,e^{-1/2}\to\infty$, так что
        $\rho_n\to\infty$ — равномерной сходимости нет.
\end{enumerate}
\end{example}

\begin{criterion}[Коши для последовательностей]
$\{f_n\}$ равномерно сходится на $X$ тогда и только тогда, когда
\[
  \forall\eps>0\ \exists N\ \forall n>N\ \forall p\in\N\ \forall x\in X:\
  \abs{f_{n+p}(x)-f_n(x)}<\eps .
\]
\end{criterion}

\begin{proof}
$(\Rightarrow)$ Если $f_n\rightrightarrows f$, то по $\tfrac\eps2$ найдётся $N$: при $n>N$
$\abs{f_n(x)-f(x)}<\tfrac\eps2$ для всех $x$; тогда
$\abs{f_{n+p}(x)-f_n(x)}\le\abs{f_{n+p}-f}+\abs{f-f_n}<\eps$. $(\Leftarrow)$ При каждом $x$
последовательность $\{f_n(x)\}$ фундаментальна, значит сходится к некоторому $f(x)$. Переходя в
неравенстве $\abs{f_{n+p}(x)-f_n(x)}<\eps$ к пределу $p\to\infty$, получаем
$\abs{f(x)-f_n(x)}\le\eps$ для всех $x$ при $n>N$, т.е. сходимость равномерна.
\end{proof}

\section{Функциональные ряды}

\begin{definition}
Ряд $\displaystyle\sum_{k=1}^\infty u_k(x)$ называют \emph{функциональным}; его частичные суммы
$S_n(x)=\sum_{k=1}^n u_k(x)$. Множество $x$, где $\{S_n(x)\}$ сходится поточечно, — \emph{область
сходимости}. Ряд сходится \emph{равномерно} на $X$, если $S_n\rightrightarrows S$ на $X$.
\end{definition}

\begin{definition}
\emph{Остаток ряда} $\displaystyle R_n(x)=S(x)-S_n(x)=\sum_{k=n+1}^\infty u_k(x)$. Удобный
критерий: ряд сходится равномерно $\iff\displaystyle\sup_{x\in X}\abs{R_n(x)}\to0$.
\end{definition}

\begin{criterion}[Коши для рядов]
$\sum u_k(x)$ сходится равномерно на $X$ тогда и только тогда, когда
\[
  \forall\eps>0\ \exists N\ \forall n>N\ \forall p\in\N\ \forall x\in X:\
  \abs[\Big]{\sum_{k=n+1}^{n+p}u_k(x)}<\eps .
\]
\end{criterion}

\begin{theorem}[признак Вейерштрасса, мажорантный]
Если $\abs{u_k(x)}\le a_k$ для всех $x\in X$ и числовой ряд $\sum a_k$ сходится, то
$\sum u_k(x)$ сходится равномерно (и абсолютно) на $X$.
\end{theorem}

\begin{proof}
По критерию Коши для $\sum a_k$: $\sum_{k=n+1}^{n+p}a_k<\eps$ при $n>N$. Тогда
$\abs[\big]{\sum_{k=n+1}^{n+p}u_k(x)}\le\sum_{k=n+1}^{n+p}\abs{u_k(x)}\le\sum_{k=n+1}^{n+p}a_k<\eps$
равномерно по $x$ — выполнен критерий Коши для $\sum u_k$.
\end{proof}

\begin{example}
$\displaystyle\sum_{n=1}^\infty\frac{\sin nx}{n^2}$ сходится равномерно на $\R$:
$\abs[\big]{\tfrac{\sin nx}{n^2}}\le\tfrac1{n^2}$, а $\sum\tfrac1{n^2}$ сходится (мажоранта).
\end{example}

\begin{remark}[признак Дирихле]
Ряд $\sum f_k(x)g_k(x)$ сходится равномерно на $X$, если частичные суммы $\sum f_k$ ограничены в
совокупности ($\abs{\sum_{k=1}^n f_k(x)}\le M$) и $g_k(x)\rightrightarrows0$ монотонно по $k$.
Пример: $\sum\dfrac{\sin nx}{n^p}$ ($p>0$) равномерно сходится на $[\delta,2\pi-\delta]$.
\end{remark}

\section{Свойства равномерно сходящихся рядов}

\begin{theorem}[непрерывность суммы]
Если все $u_k$ непрерывны на $X$ и $\sum u_k\rightrightarrows S$, то $S$ непрерывна на $X$.
\end{theorem}

\begin{proof}[Метод «трёх $\eps$»]
Фиксируем $x_0\in X$ и $\eps>0$. Так как $S=S_n+R_n$,
\[
  \abs{S(x)-S(x_0)}\le\abs{S_n(x)-S_n(x_0)}+\abs{R_n(x)}+\abs{R_n(x_0)} .
\]
В силу равномерной сходимости выберем $n$, при котором $\abs{R_n}<\tfrac\eps3$ всюду. При этом $n$
сумма $S_n$ непрерывна (конечная сумма непрерывных), поэтому $\exists\delta>0$: при
$\abs{x-x_0}<\delta$ имеем $\abs{S_n(x)-S_n(x_0)}<\tfrac\eps3$. Итог:
$\abs{S(x)-S(x_0)}<\tfrac\eps3+\tfrac\eps3+\tfrac\eps3=\eps$.
\end{proof}

\begin{theorem}[почленное интегрирование]
Если $u_k$ непрерывны на $[a,b]$ и $\sum u_k\rightrightarrows S$, то
\[
  \int_a^b S(x)\dx=\int_a^b\sum_{k=1}^\infty u_k(x)\dx=\sum_{k=1}^\infty\int_a^b u_k(x)\dx .
\]
\end{theorem}

\begin{proof}
$S$ непрерывна, значит интегрируема. Оценим разность:
\[
  \abs[\Big]{\int_a^b S-\int_a^b S_n}=\abs[\Big]{\int_a^b R_n}\le\int_a^b\abs{R_n}
   \le\int_a^b\frac{\eps}{b-a}\dx=\eps
\]
при $n>N$ (равномерность). Значит $\int_a^b S_n\to\int_a^b S$, т.е. интеграл суммы равен сумме
интегралов.
\end{proof}

\begin{theorem}[почленное дифференцирование]
Пусть $u_k\in C^1[a,b]$, ряд $\sum u_k'(x)\rightrightarrows\sigma(x)$ равномерно на $[a,b]$, и
$\sum u_k(x_0)$ сходится хотя бы в одной точке $x_0$. Тогда $\sum u_k\rightrightarrows S$ на
$[a,b]$, $S$ дифференцируема и $S'(x)=\sigma(x)=\sum u_k'(x)$.
\end{theorem}

\begin{proof}
По теореме о почленном интегрировании равномерно сходящегося ряда $\sum u_k'$:
$\int_{x_0}^x\sigma(t)\dt=\sum_k\bigl(u_k(x)-u_k(x_0)\bigr)$, т.е.
$\sum u_k(x)=\int_{x_0}^x\sigma+\sum u_k(x_0)$. Правая часть дифференцируема с производной
$\sigma(x)$ (по теореме Барроу), значит $S'=\sigma$.
\end{proof}

\begin{remark}
Условие равномерной сходимости \emph{ряда производных} существенно: его отсутствие делает
почленное дифференцирование недопустимым. Так, для $f_n(x)=\tfrac{\sin nx}{\sqrt n}$ имеем
$f_n\rightrightarrows0$, но $f_n'(x)=\sqrt n\cos nx$ не стремится к нулю.
\end{remark}
